% !TeX root = ../incremental_SS_Translation.tex

\section{Uttaratantra 64: Preservation of Health}% svastharakṣanīya

\section{Literature}

Meulenbeld offered an annotated overview of this chapter, which is called 
\emph{svasthavṛtta} in the vulgate.\fvolcite{IA}[330--331]{meul-hist} There, 
he translates \emph{svasthavṛtta} as ``preservation of health'' but elsewhere 
as ``maintenance of health''\fvolcite{IA}[301]{meul-hist} and 
``hygiene.''\fvolcite{IB}[738]{meul-hist} His only reference to secondary 
literature on this topic is to an annotated German translation by J. Laping 
1984.\footnote{\cite{lapi-1984}, cited by \volcite{IB}[431]{meul-hist}.} His 
footnotes to the chapter also indicate 
many parallels with \emph{Carakasaṃhitā, Sūtrasthāna}, chapter six.

On the term \emph{svastha}, Meulenbeld cites W. Halbfass (1991): 250-252; M. 
Hara (1995) and P.V. Sharma (1985e: 52-53).\fvolcite{IB}[10]{meul-hist}

In primary sources, see the definition of \emph{svasthā} in 
\emph{Suśrutasaṃhitā, Sūtrasthāna} 15.41--44; \emph{svasthahita} in 
\emph{Aṣṭāṅgahṛdaya, Sūtrasthāna} 1.16cd; ...

Comparable chapters in other medical texts include the \emph{svasthacatuṣka} 
in \emph{Carakasaṃhitā, Sūtrasthāna} (see below) and a chapter on 
\emph{svasthavṛtti} in Cakrapāṇidatta's \emph{Cakradatta}.

Meulenbeld also refers to \emph{svasthacatuṣka}, a group of four chapters 
(5--8) in the \emph{Carakasaṃhitā, Sūtrasthāna}, chapters 5–8, also called 
\emph{svasthavṛttacatuṣka} or \emph{svāsthyavṛttikacatuṣka} in various 
colophons of the text and commentaries. He summarises these 
chapters,\fvolcite{IA}[13--14]{meul-hist} and their names are 
\emph{mātrāśitīya}, daily regime (cf. Bhela Sū.8; A.h.Sū.8; A.s.Sū.11. Cf. 
Su.Ci.24); \emph{tasyāśitīya}, seasonal regimen (\emph{ṛtucaryā}) 
(Meulenbeld cites many references in primary and secondary 
literature\fvolcite{IB}[13]{meul-hist}); \emph{navegāndhāraṇīya}, the 
unwholesome effects of suppressing natural urges, etc. (Cf. Bhela Sū.6; 
Su.Ci.24); and \emph{indriyopakramaṇīya}, senses, mind and rules of conduct 
(Cf. Bhela Sū.7; Su.Ci.24. See on Ca.Sū.8: S.Ch. Vidyabhusana (1971): 11--12).

Meulenbeld mentions that a ``Sadānanda wrote a commentary, called \emph{Auṣadhavivṛti}, on the chapters of the
\emph{Carakasaṃhitā} containing rules for the preservation of health 
(svasthavṛtta).\fvolcite{IA}[196]{meul-hist} His footnote on this work refers to 
the NCC (VI, 396–397) and the following bibliographic reference: 
\emph{pañcatantram \ldots\  
carakasūtrasthānasvasthavṛttacatuṣkākhyacaturadhyāyāḥ\ldots\
sadānandaśāstrikṛtauṣadhavivṛtiyutayā saṃvalitaṃ}, Mercantile Press, Lahore 
1926 [IO.San.D.554].\fvolcite{IB}[302]{meul-hist}

In the vulgate, this chapter is included in a section called, ``the chapters that 
are an embellishment of the text" (\emph{tantrabhūṣaṇādhyāya}). However, in 
the Nepalese version, it is the eighth part of tenfold section on treatment for the 
body (\emph{kāyacikitsita}).\footnote{See witness K, folio 209r, ll. 4--5 
(\dev{mūtradoṣo mūtrāghāto yonyamānūṣa eva ca/  
apasmāronmādakañcaiva rasabhedastathaiva ca// 
svastharakṣāvidhāṇañca tantrayuktiśca doṣabhit/
ityebhirdaśakaṃ proktaṃ punaḥ kāyacikitsite//}).}

This chapter is cast as a discussion between Dhanvantarī and the king, who is not named but is likely to be Divodāsa. 

\section{Translation}
% K f. 205v
% H f. 406

\begin{translation}    
  
\item [1] 
Now we shall explain how to preserve a person's health.%
% the upkeep of a healthy person.% rakṣaṇīya - maintenance of? prophylaxis for a healthy person?
	\footnote{The term \emph{svastharakṣaṇīya} more literally means, ``preserve a healthy person.'' \Dalhana{6.64.1}{808} noted this reading, which is also what we see in the Nepalese version, and explained it in these terms.
    
    The word \emph{rakṣaṇīya} (preservation, maintenance, upkeep, etc.) is replaced in the vulgate \citep[808]{vulgate} by \emph{vṛtta}, which, when compounded with \emph{svastha}, denotes the “regime for a healthy person.” However, \emph{rakṣaṇa} is retained later in the same chapter of the vulgate.} 

\item [2] 
One whose humours, digestive fire, and the functioning of bodily constituents and waste products are in equilibrium; whose self, senses, and mind are serene is called healthy.\footnote{This verse is absent here in the vulgate but occurs at the end of \emph{Sūtrasthāna} 15 in both the Nepalese version and vulgate. The vulgate \citep[808]{vulgate} supplements its omission with the statement, ``Such a healthy person has been described in the \emph{Sūtrasthāna}'' (\emph{sūtrasthāne samuddiṣṭaḥ svastho bhavati yādṛśaḥ}). \Dalhana{1.15.41}{75--76} discussed several points of view
regarding the meaning of equilibrium in these different 
contexts.}\q{More on Ḍalhaṇa and the meaning of sama?}

\item [3] 
My dear, the preservation of a healthy person is the supreme benefit of medicine. The topic of his regime and preservation, which was mentioned by me at the beginning, was outlined in brief and will here be explained in detail.
% The interlocutors are king (divodāsa) and Dhanvantari.See the end of the chapter.
% something has gone wrong with artham in the Nepalese manuscripts (which has arthāṃ and arthān). The participle (uktaṃ) and verb (vakṣyate) require a singluar. And artham can be neuter. The syntax of the vulgate is wrong because the subject is plural (arthāḥ) and the verb singular (a patch made inadequate by metric causa?)
% ḍalhaṇa: āditaḥ anāgatabādhe

\item [4] 
In each season, when specific humours are aggravated in people, a wise physician should prescribe the appropriate tastes for that season.

\item [5] % varṣa season
You should know that when people's digestive fire is weak because of the dampness of their bodies during the rainy months, wind and other humours become aggravated due to their excitement.\footnote{%
    The vulgate notes that the readings \emph{saṃgharṣāt} and \emph{saṃharṣāt} occur in an unspecified number of witnesses (\cite[809, n.\,2]{vulgate}). The reading \emph{saṃgharṣāt} (``friction'') makes better sense in the context of the humours, but is not supported by the Nepalese manuscripts.}
% Ḍalhaṇa appears to understand mārutādayaḥ as referring to the doṣas.
% have I over-translated khalu?

\item [6--8ab] 
Therefore, in order to remove moisture and subdue the humours,
the king should be given solid food with astringent, bitter, and \se{kaṭu}{spicy} tastes, which is neither too oily nor too dry, and which is hot and \se
{dīpana}{stimulates digestion}.\footnote{%
        The vulgate reads \emph{doṣasaṃ\textbf{hara}ṇāya} (``for removing the humours'') instead of \emph{doṣasaṃ\textbf{dhāra}ṇāya} \citep[809]{vulgate}. \CS\ \Ca{2.6.6}{219--220} uses the word \emph{saṃdhāraṇa} in the context of restraining the urges to urinate and so on. In the current context, \emph{doṣasaṃdhāraṇa} refers to subduing the humours that have been aggravated by weak digestive fire during the rainy season (mentioned in the previous verse).}  
He should drink the water which was mentioned before,\footnote{%
        \Dalhana{6.64.8}{809} understood “the water which was mentioned before” to be monsoon rainwater, which was described in the section on what to drink (\dev{yajjalaṃ coktamādita iti ‘varṣāsvantarikṣaṃ’ ityādinā pānīyavarge kathitam}). This is the section that starts at \SS\ \Su{1.45.3–46}{196--200}.} %
or water which has been treated with hot iron\q{%
        Look for references to treating water with hot iron.} %
and mixed with honey.%
%

% See Roots Ed 3 p 228 for the tastes and the passage in the Aṣṭāṅgahṛdaya on the tastes
% 

\item [8cd--9]
On a rainy and windy day that is very cold and wet, % atyarthaśītāmbusaṃkule: literally: affected by excessively cold water
herbal medicines develop acidity because of being fresh.%
	\footnote{Ḍalhaṇa \citep[809]{vulgate}  glossed \emph{vidāham} (“burning”) as \emph{amlapāka} (“sour digestion”), suggesting that fresh herbs may cause acidic irritation in the stomach when taken on a rainy, windy, cold, damp day.}
	%
And because of that, a sensible person should not exercise excessively.

\item [10] % not in the vulgate
When it is not very cold, aged barley, wheat and rice with spicy broths and soups, as well as appetizing meats from dry land, 
should be eaten.\footnote{``The meats of dry land'' are those of \sepl{jaṅgala}{arid land}; see \cite[fig.\,5, p.\,52]{zimm-1999} for a list of the \SS's meats in this class.} 
% emended to bhojyā based on Caraka Sū 6.38c (note in the edition) 
% the ca in pāda 4 suggests that jāṅgala is wild meat/deer/game rather than an adjective. Elsewhere in Ḍalhaṇa I see eṇādayo jāṅgalamṛgamāṃseṣu (an adjective here but obviously refers to game).

 % start here next time
\item [11] % not in the vulgate
One should drink a little mead or sweet wine, %
	\footnote{On \emph{mādhvīka}, a kind of honey liquour, see McHugh 2021: 46. On \emph{sīdhu} (“sweet wine”), Ḍalhaṇa \citep[778]{vulgate} says it is \emph{gauḍika}, an alcoholic drink produced from sugar or grain (\emph{gauḍikaṃ sīdhu guḍakṛtaṃ sasyakam}). McHugh (2021: 45–46), who translates \emph{sīdhu} as “sugarcane wine,” says that it is primarily made from sugarcane but also other sweet substances, like jaggery, herbs and mahua flowers.}
%
and also rainwater, lake water, or well water, hot or cold and mixed with honey. 

\item[12] % ab not in the vulgate
At this time, % atra or 'In this climate'?
one should avoid the easterly wind, exposure to the heat, buttered barley water, river water, sleep during the day, frost, and sex.

\item[13--14ab] % not in the vulgate
One should engage in body rubs, massages, bathing, perfumes, garlands and smoking. To avoid humidity from the ground, one should have a heavy cloak and sleep in a house's central room, which is cool, has a fire, and is free from draughts.
%
% uṣma heat, hot vapour, steam. Is "humidity" a better translation  here as I think this still concerns the varṣa season?
% cf: Caraka1.6.40: pragharṣodvartanasnānagandhamālyaparo bhavet | laghuśuddhāmbaraḥ sthānaṃ bhajed akledi vārṣikam || [during the rainy season]
% cf: AS.Sū.4.48	pragharṣodvartanasnānadhūmagandhāgarupriyaḥ | 
%				yāyāt kareṇumukhyābhiścitrasragvastrabhūṣitaḥ || [varṣa season]
% Jayaratha on Tantrāloka: madhyegṛhamiti--gṛhamadhye brahmasthāne--ity arthaḥ
% madhyamaveśmani might denote 'in an inner/central chamber of the house.'
% pragharṣa = gharṣaṇa (Āyurvedīya Kośa), and gharṣaṇa = mardana. mardana = aṅgamardana
% udvartanaṃ = gātramardanaṃ (Āyurvedīya Kośa); 1. tvagavagharṣaṇa 2. snānāt pūrvaṃ tvagavagharṣaṇārtham upayujyamānaṃ cūrṇam āmalakyādi. 3. pādābhyām aṅgāvayavānāṃ pīḍanaṃ yena vyāyāmajanitaḥ śramaḥ śāmyati.

\item[14cd] % not in the vulgate
And if one is injured, one should avoid morbid conditions (\emph{doṣa}), such as moisture, and be protected by an umbrella.

\item[15ab]
One also should travel by means of female elephants and adorn oneself with the excellent aloe-wood.%
	\footnote{Ḍalhaṇa \citep[809]{vulgate} remarks that one should travel on female elephants to avoid mud and water, and one should wear aloe-wood because it wards off the cold (\emph{nāgavadhūbhir iti hastinīgamanaṃ ca kardamāmbunoḥ parihārārtham| praśastāgurubhūṣita iti aguruṇaḥ, śaityanivārakatvāt}).} %
Thus, [it has been taught].%
	\footnote{The \emph{iti} after 15ab appears to follow from verse 3. In K, the \emph{iti} is followed by a rubricated section marker.}\\

%
% Then the vulgate has:
% divāsvapnam ajīrṇaṃ ca varjayet tatra yatnataḥ | 
% sevyāḥ śaradi yatnena kaṣāyasvādutiktakāḥ ||
% kṣīrekṣuvikṛtikṣaudraśālimudgādijāṅgalāḥ |
% śvetasrajaś candrapādāḥ pradoṣe laghu cāmbaram ||
% salilaṃ ca prasannatvāt sarvam eva tadā hitam
% saraḥsv āplavanaṃ caiva kamalotpalaśāliṣu ||

\item[15cd--17ab]
Homage to the sun, who has spread forth its light along the bright pathways of the quarters, when the sky's covering of dense blue is spread over the extent of the earth, and who is exalted by being born at the sight of Indra’s thunderbolt. The splendour of the sun advances, smothering cloud after cloud.
% I am assuming the locative is being used for the dative. Oberlies observes that namaskṛtya can take the locative in Epics (p 353), but also that the locative and dative are exchangeable. 
% not sure about °vaguṇṭhite vyomni. I've translated avaguṇṭhite more like a noun. 
%\skt{labdha-prasara}    mf(\skt{ā})n. that which has obtained free scope, moving at liberty, unimpeded Mudr. But does it  refer to the sun rising (i.e., when the sun has attained its rising) or the spreading forth of its light?

\item[17cd--18ab]
 Bile that is accumulated during the rains in autumn is aggravated by the heat. For this reason, when the clouds are dispersing, all people should undertake practices that reduce bile. 

\item[18cd--19ab]
During autumn, when one has an appetite, one should eat food that is cool, astringent and sweet, somewhat light and  bitter, as well as mixed with ghee and containing game.
% can food be both astringent and sweet or should we understand the compound kaṣāyamadhuram as astringent or sweet?
% I have read  īṣal in compound with laghu. 
% Is satiktakam bitter or pungent?
% jāṅgala appears to range from deer (jaṅghāla) to birds and burrowing animals (Ḍalhaṇa 1.46.53: jāṅgalaśabdena jaṅghālaviṣkirapratudaparṇamṛgabileśayā iti trayodaśa bhedāḥ ṣaṭsvevāntarbhūtāḥ)

\item[19cd--20ab]
If not indulging excessively in sex, one should enjoy eating deer, spotted antelopes, quails, goats, black antelopes, and  partridges, along with mung beans, aged rice, and barley.

\item[20cd--21ab]
One should avoid burning, acrid, spicy and heating foods, sleeping during the day, and heavy foods, as well as vegetables and meats that are heavy, and exposure to frost.

% It seems strange to have divāśvapna in the middle of a compound of types of food, so I suspect that at least one anusvāra has dropped out, more likely two (unless I've misunderstood gurūṇi. Note my emendations. The atha suggests at least two compounds were intended here (instead of the one in K and H).
% Cf. Aṣṭāṅgahṛdaya, nidānasthāna 13.2526
% vyādhikarmopavāsādikṣīṇasya bhajato drutam | 
% atimātram athānyasya gurvamlasnigdhaśītalam | 
% lavaṇakṣāratīkṣṇoṣṇaśākāmbu svapnajāgaram |
% mṛdgrāmyamāṃsavallūram ajīrṇaśramamaithunam ||

\item[21cd--22ab] % not in the vulgate
O King, in autumn, the wise person should carefully follow this rule:
One should consume grapes, sugarcane and milk in that season (\emph{tatra}).
% confirmation that a king is being spoken to here.

\item[22cd--23ab]
When bile has accumulated during the rains, one should eliminate it with purgatives, or by administering bitter medicated ghee, and also by bloodletting from the veins.
% I assume that tiktaḥ and sarpis need to be read in compound with prayoga (and have emended tikta accordingly). 

\item[23cd--24] % In the vulgate marked by round brackets
Water should be tasty and cool, pure, clear like a flawless crystal, cleansed by the rays of the autumn sun, and detoxified by the rising star Agastya. Because of its clearness, all such water is then beneficial.

\item[25--26ab] %  In the vulgate marked by round brackets
And in autumn seasons, one should enjoy the breeze in lotus, lily and rice fields; or the sandalwood with fragrant camphor; clean light clothes; autumnal garlands; and within reason drinking sweet wine.
%
% pavanam ... kamalotpalaśāliṣu: the breeze blowing among lotuses, lilies and rice stalks?
% yuktitaḥ – within reason?

\item[26cd--27ab] % 26cd not in vulgate
One should do everything that pacifies bile: enjoying the moonlight, a circle of friends, pleasant comforts, and sweet speech. Thus it has been said.

\item[27cd] 
Winter is cold, dry, with a mild sun, and blustery.

\item[28--29] 
One should consume the following food and drinks: sesame seeds, green gram, vegetables and curds, as well as sugarcane products, and rices that are fragrant and also fresh, and meats of carnivorous and marsh-dwelling animals, carnivorous burrow-dwelling creatures, aquatic animals, birds and amphibious creatures.  
% Emended māṣāṃ śākāni (K) to māṣāñ śākāni
% emended vikṛtīñ (K: pc) to vikṛtīḥ (vulgate).
% I don't understand the syntax here as the genitive plurals appear to depend on māṃsāni, but māṃsāni is the final member of a compound. Maybe the original reading was something like prasahānāṃ ca māṃsāni.
% I've translated kravyāda as an adjective. Not sure if that's right, but we also have prasaha as carnivorous animals.

\item[30] 
If one desires nourishment at the approach of winter,  one may partake of clear alcoholic drinks and whatever gives strength, according to one’s inclination. 

\item[31--32ab] % 31cd and 32ab are absent in the vulgate
Bathing in moderately warm or comfortably warm oil is recommended. Alternatively, one should engage in oil massage including the head, perform vigorous exercise, bathe in comfortably warm water and wash with only warm water. 
% I've understood the vā is expressing an important alternative here to bathing in oil

\item[32cd--34ab] 
One whose body has been smeared with saffron and aloe-wood % kuṃkumāgurudigdhāṃgo. Is this referring to some sort of paste made with saffron and aloe-wood?
should rest % śayīta
on a large %mahati
bed, covered with wool and silk fabrics % aurṇṇakauśeyasamvītaśayane
and overspread with a dyed blanket. % kuthakāstṛte. Ganguly note: Kuthaka* is a painted blanket. PV Sharma, variegated blanket (translating Caraka 1.6.15cd)
It has a vessel for coals, % sāṃgārayāne. Is there any secondary literature on Indian beds in the classical era that might illuminate this? PV Sharma takes it with garbhagṛhodara: with mobile containers of burning charcoals. But even in the very different version of the vulgate's verse, this seems unlikely (word order of the vulgate still suggests it qualifies śayane). Srikantha Murthy translates AS.Sū.4.17 (sāṅgārayānāṃ śayyāṃ) as 'equipped with an oven' but still construes it with house, even though the syntax indicates it qualifies a bed. It seems more likely to me that an āṃgārayāna must be a vessel with hot coals in it put under or near the bed to keep it warm (or may actually be built into the bed, like some Tibetan beds). Maybe "brazier with coals" is the right word?
is filled with the strongly fragrant musk of incense, % dhūpāmodamadotkaṭau: mada; musk. āmoda; strong fragrance. utkaṭa; richly endowed with, abounding in. 
 and is in an inner chamber of the house, % gṛhagahvare
free from drafts, % nivāte
And the bed should also be % cāpi śayane. I am tempted to emend to vāpi. As the Sanskrit stands, it seems that two beds are being referred to here, but I think the repitition of śayane is a result of poor redaction.
very wide % suvistīrṇṇe
and beautiful. % manorame

\item[34cd--35] % not in the vulgate
In it, % tatra
devoted % priyā
women who are well-adorned % nāryaḥ svalaṃkṛtāḥ
with their necklaces removed % apanītahārās 
and intensely passionate, % madotkaṭāḥ
should delight him, % rāmayeyur
at the proper time, % yathākālaṃ
even by force, % balād api
while he is being attended to %sevyamānaṃ
with gentle touches % mṛduśparśair
and aggressive embraces. % nirdayair upagūhanaiḥ: nirdaya - merciless, cruel, violent, excessive
% Cakrapāṇi quotes part of this verse: yaduktaṃ- “tatrāpanītahārāś ca priyā nāryaḥ svalaṅkṛtāḥ | ramayeyur yathākāmaṃ balād api madotkaṭāḥ | This justifies many of the emendations I have made.

\item[36] 
This procedure % eṣa vidhi
that has been taught should be done during the dewy season. % śiśira: the cool or dewy season (comprising two months, Māgha and Phālguna, or from about the middle of January to that of March)
Some physicians consider the remaining part of winter to be the dewy season. 

\item[37] % not in the vulgate
Since the cold [of the dewy season] is produced by clouds and wind and the dry caused by the warm seasons,%
	\footnote{The term \emph{ādāna} has been translated as ``warm seasons'' because it is generally understood to be the warm half of the year or the three warm seasons during which the sun takes (ā + ⁄{dā}) strength from living beings. The term appears in \emph{Carakasaṃhitā, Sūtrasthāna} 6.4, where it is defined as the three seasons beginning with the cool season (\emph{śiśira}) and ending with summer (\emph{griṣma}). The  \emph{Āyurevedīyakośa} (s.v.) defines \emph{ādāna} as, ``the three seasons of spring, summer and the rains, during which the sun takes strength from living beings''  (\emph{vasaṃtagrīṣmavarṣās traya ṛtavaḥ ādānakālaḥ asmin sūryaḥ prāṇināṃ balam ādatte}).}
	%
one who desires a long life should adopt a regime that is more heating and moisturising. 	

% ādāna: P.V. Sharma translates raukṣyam ādānajaṃ in Caraka (1.6.19cd, below) "roughness due to (beginning of) ādāna. Big help. 
% According to the Āyurevedīyakośa (s.v.): vasaṃtagrīṣmavarṣās traya ṛtavaḥ ādānakālaḥ asmin sūryaḥ prāṇināṃ balam ādatte. So, it is the three seasons of Spring, Summer and Rains, during which the sun takes strength from living beings. 
% Āyurvedīyakośa also says: kālādānaḥ ādadāti kṣapayati pṛthivyāḥ saumyāṃśaṃ prāṇināṃ ca balam ity ādānam
% Cf. Caraka
% 1.6.19ab	| hemantaśiśirau tulyau śiśire 'lpaṃ viśeṣaṇam |
% 1.6.19cd	| raukṣyam ādānajaṃ śītaṃ meghamārutavarṣajam ||
% 1.6.20ab	| tasmād dhaimantikaḥ sarvaḥ śiśire vidhiriṣyate |
% 1.6.20cd	| nivātam uṣṇaṃ tv adhikaṃ śiśire gṛham āśrayet ||
%
% Cf. AS.Sū.4.20	
% raukṣyaṃ ca+ādānajaṃ tasmāt kāryaḥ pūrvo 'dhikaṃ vidhiḥ | 
% vasante dakṣiṇo vāyurātāmrakiraṇo raviḥ ||
% Srikantha Murthy translates it as the Ādāna kāla. 
%
% In summarising chapter 6 of Caraka Sūtrasthāna, Meulenbeld writes:
% The subjects are: the importance of seasonal regimen (6.3); the year is divided into six seasons; the period in which the sun courses northwards (udagayana), known as ādāna, consists of the seasons beginning with śiśira (the cool season) and ending with grīṣma (summer); the period in which the sun courses southwards, called visarga, consists of the seasons beginning with varṣāḥ (the rainy season) and ending with hemanta (winter) (6.4); the main characteristics of visarga and ādāna (6.5-8); ...
% In footnote 119 on the main characteristics of visarga and ādāna, Meulenbeld says: During visarga the rays of the moon replenish the world, thus restoring it after a period of ādāna by the sun. 
% When commenting on the Aṣṭāṅgasaṅgraha, Sūtrasthāna, ch. 4, Meulenbeld says: The first half of the year is called udagayana or ādāna, the second half dakṣiṇāyana or visarga. The āgneya character of ādāna and the saumya character of visarga are described (in prose, concluded by a verse) (4.5–7).

\item[38] % not in the vulgate
During this time (\emph{tatra}), one should avoid mixed barley grout drinks,%
	\footnote{The term \emph{mantha}, which literally means ``stirred,'' is defined in Monier-Williams dictionary (1899: s.v.) as '`mixed beverage, usually parched barley-meal stirred round in milk.'' In two instance, Ḍalhaṇa \citep[811 and 779]{vulgate} glosses \emph{mantha} as barley groats prepared with water and ghee (\emph{manthān jalaghṛtāktasaktūn}), and barley grouts covered with cold water along with ghee (\emph{manthaḥ saghṛtaśītalajalaplutāḥ saktavaḥ}).} %
%
frosts, windy weather, dry foods, wearing dirty clothes, and bathing. %
%
Thus, it has been said.

\item[39--41] % not in the vulgate
While the rays of the blazing sun % dīptabhāskararaśmiṣu
are melting % nirhāra = nirhāraṇa (MW): removing
the layers of frost, % tuṣārapaṭa
and the days are lengthening and nights diminishing, % dineṣu jṛmbhamāṇeṣu hīymānāsu rātriṣu
%
when the wondrous  % vicitrāsu: variegated, diverse, but here maybe wondrous given the metaphors that follow
rows of forests and groves % vanopavanarājiṣu
are displaying [themselves], % darśayatsu 
with colours variegated % śabalaiḥ
by the loud laughter of flowers, % puṣpāṭṭahāsa ?
radiant with trembling young shoots, % calat-kisalaya-ujjvalaiḥ
and sounds % niḥsvanaiḥ
resonating with cuckoos % kokilodghuṣṭa
and the humming of intoxicated bees, % samadabhramarodgīta
%
phlegm becomes aggravated % kaphaḥ prakupyati
and will produce % sisṛkṣati -desiderative? Whitney (948b, 1040a) says the desiderative is sometimes found where the future tense is expected.
the diseases as described. % yathoddiṣṭān ... gadān

\item[42] 
Cold phlegm that has accumulated in people during winter because of the cold is dissolved % viṣyaṇṇaṃ (vi-ṣyand, literally,  to dissolve, melt)
in spring due to the heat and causes various diseases.

\item[43] 
Because of this, one should avoid water, sweet and oily food, sleep during the day, and heavy foods, and one should also undergo treatments, such as emesis.





\end{translation}    
